\documentclass[12pt]{article}

% House style preamble
\input{.house-style/preamble.tex}

% Update PDF metadata
\hypersetup{
  pdftitle={Truth-Tracking as a Homeostatic Property Cluster: What LLMs Participate In}
}

\title{Truth-Tracking as a Homeostatic Property Cluster:\\ What LLMs Participate In}
\author{Brett Reynolds \orcidlink{0000-0003-0073-7195}%
\thanks{Contact: \href{mailto:brett.reynolds@humber.ca}{brett.reynolds@humber.ca}}\\
Humber Polytechnic \& University of Toronto}
\date{\today}

\begin{document}
\maketitle

\begin{abstract}
% TODO: scaffold in docs / NOTES.md. The claim is NOT "truth is an HPC kind" but
% "truth-tracking -- the family of representational successes picked out by the
% predicate \mention{true} -- is stabilised by an HPC-like cluster of mechanisms,
% and LLMs are the case that makes a cluster-theoretic account necessary."
TODO: Write abstract here.
\end{abstract}

% ============================================================================
% STRUCTURE NOTE. This skeleton encodes the design decisions in DECISIONS.md and
% the full brief in NOTES.md. Section 1 must do the load-bearing non-relabeling
% work (answer "what does HPC explain that reliabilism / coherentism / pragmatism
% / correspondence don't?") -- it is NOT a generic scene-setting intro.
% Stubs only below; do not draft prose without reading sources (source-grounding LAW).
% ============================================================================

\section{Introduction: a binary presupposition}\label{sec:intro}

Do the representations inside a large language model connect to the world? The
question is usually put as a yes/no. \textcite{bender2020} say no: a system trained on
form alone, like an octopus eavesdropping on telegraph traffic, learns the statistics
of language without touching what language is about. Others reply that such systems
meet at least some conditions on reference, or that sensory grounding isn't required
for thought in the first place
\citep{coelho_mollo_milliere_2026_vector_grounding_problem,chalmers_2023_thought_sensory_grounding_llms}.
The camps disagree on the verdict but share the question's shape.
% TODO: after close reading, fold Pepp 2025, Grindrod 2024, Mallory 2026, and
% Mandelkern & Linzen (needs a bib entry) into the rebuttal-side citation here.

Both sides assume that connecting to the world is a single capacity, something a system
either has or lacks. That assumption is the mistake. Truth-tracking isn't one capacity;
it's a cluster of distinct mechanisms that ordinarily operate together: perceptual
contact, testimonial inheritance, inferential coherence, measurement and intervention,
error-correction, and practical success. A system can take part in some of these while
sitting out the rest, which is just what a language model does.

The paper's claim is that truth-tracking, the family of representational successes we
pick out with the predicate \mention{true}, is stabilised by a homeostatic property
cluster (HPC) in Boyd's sense \citep{boyd1991,boyd1999}. What makes the mechanisms a
cluster and not a list is that they discipline one another: perception corrects
testimony, measurement corrects coherence, and practical failure corrects them all.
That mutual disciplining holds the kind together, and it's what makes the cluster's
instances project.

This develops Boyd's programme rather than rivalling it.
\textcite{boyd_2019_rethinking_natural_kinds_reference_truth} already treats reference
and truth as special cases of accommodation between our signalling practices and the
causal structures they answer to. The cluster view of truth-tracking applies that
picture one level up. Boyd's accommodationism is the background; the contribution here
is the diagnostic the cluster structure supports once a system enters some accommodation
practices but not the others.

That diagnostic has to be earned, because the cluster view invites an obvious complaint.
Reliabilism already traces true belief to reliable mechanisms; coherentism names
inferential fit; pragmatism names practical success; social epistemology names
testimony. Listing these and calling the bundle \enquote{homeostatic} can look like
relabeling, four familiar theories in one new coat. If that's all the view does, it adds
nothing.

It does more, in three ways no single-mechanism theory can match. First and most
sharply, the cluster structure predicts a characteristic failure signature: pull the
stabilisers apart, and the account says in advance what should go wrong. Second, it
makes conflicts between mechanisms adjudicable, by treating the degree to which they
co-vary as itself evidence. Third, it predicts where truth-relevant mechanisms should
co-vary tightly across domains and where they should come apart.

The language model is where the first prediction pays off. A text-trained system
inherits testimony and strong sensitivity to textual coherence, but it has no perceptual contact,
runs no measurements, and gets little correction from acting in the world. The cluster
view says such a system should produce claims that are confident, fluent, internally
coherent, faithful to corpus consensus, and yet unanchored to fact. That profile has a
name: hallucination.

On this account hallucination is the predicted output of a partial cluster, not a defect
to be patched piecemeal. It marks the distance between the formal linguistic competence
these models plainly have and the functional competence that world-contact underwrites
\citep{mahowald2024}.

The argument runs as follows. Section~\ref{sec:target} fixes the target: the members of
the cluster are representational successes (beliefs, utterances, measurements, models,
inferential outputs), not truth as an abstract object. Section~\ref{sec:cluster} sets out
the stabilising mechanisms and their Boydian lineage. Section~\ref{sec:projectibility}
says why the cluster projects, locating projectibility in goal achievement under
counterfactual perturbation rather than in survival alone. Section~\ref{sec:adjudication}
makes good on the adjudication and co-variation claims, the core of the reply to
relabeling. Section~\ref{sec:objections} takes objections. Section~\ref{sec:llm} develops
the language-model profile mechanism by mechanism. Section~\ref{sec:conclusion}
concludes.

\section{The target: what \texorpdfstring{\mention{true}}{true} picks out}\label{sec:target}
% Pre-empt the category-error objection. Members of the cluster are representational
% successes (beliefs, utterances, inscriptions, measurements, models, inferential
% outputs), not truth-as-abstract-object. The mechanisms stabilise the EXTENSION of
% truth-ascription and explain its projectibility; they do not individually define truth.
% This blocks deflationism, pragmatist reduction, and crude evolutionary debunking.
TODO.

\section{Truth-tracking as a homeostatic property cluster}\label{sec:cluster}

The stabilisers of truth-tracking form a causal network in Khalidi's sense: each
mechanism is a node, the links between them are causal, and the kind is the densely
connected region where the nodes converge
\citep{khalidi2013,khalidi_2018_natural_kinds_nodes_causal_networks}.
Boyd's homeostatic property cluster is the special case in which those links are actively
maintained, where the mechanisms correct one another so that the properties co-occur far
more reliably than chance \citep{boyd1991,boyd1999}. Truth-tracking has such a
homeostatic core, and the connections thin toward its edges.

That picture does two jobs. The homeostatic core, everyday perception and basic
coordination, is where mutual correction is tight, and the coupling there is what
predicts the language model's failure signature: remove the perceptual and verificational
nodes and the maintained cluster comes apart in a specific way
(Section~\ref{sec:llm}).

It places the loose cases. Toward the edges, in theoretical physics, ethics, and
large-number mathematics, the nodes are real but sparsely linked. Reading the kind as a
network doesn't by itself fix where its boundaries fall: \textcite{onishi_serpico_2022}
show that carving a kind from a causal graph stays relative to the chosen level of
abstraction, the interest-relativity \textcite{craver_2009} found in mechanism talk. So
the paper fixes the purpose: truth-tracking is individuated by representing the shared
world well enough for goal achievement under perturbation
(Section~\ref{sec:projectibility}). Relative to that purpose, its connections can thin at
the edges without the kind dissolving \citep{khalidi2013}.
% TODO: enumerate the stabilising mechanisms as nodes (perceptual covariance, testimonial
% inheritance, inferential coherence, intervention/measurement, error-correction,
% practical success, counterfactual robustness); show no single mechanism is necessary;
% develop the Boydian accommodation lineage (signalling practice <-> causal structure).

\section{Why the cluster is projectible}\label{sec:projectibility}
% Projectibility story: goal achievement under counterfactual perturbation, with
% survival/reproduction as a LIMITING biological case, not the whole explanation.
% Pre-empt Hoffman-style "fitness beats truth"; cite Martinez (usefulness drives
% representations to truth) -- relation must be modelled, not asserted. In science,
% testimony, engineering, medicine, law, and LLM eval the homeostasis is institutional,
% technological, or task-relative.
TODO.

\section{What the kind-structure does that component theories don't}\label{sec:adjudication}
% SPINE ONLY (drafted 2026-05-24). Topic sentences and argument skeleton are real;
% the bracketed TODO markers flag what still needs sources and worked examples. This
% section discharges promises (2) and (3) from Section 1; promise (1), the failure
% signature, is discharged in Section~\ref{sec:llm}.

This section makes good on the claim from Section~\ref{sec:intro} that the cluster view
earns its keep. The test isn't whether \enquote{truth depends on many mechanisms} sounds
plausible; it's whether the kind-structure answers questions that reliabilism,
coherentism, pragmatism, and correspondence can't pose from inside their own resources.
Read as a causal network (Section~\ref{sec:cluster}), the kind has structure a
single-mechanism theory cannot represent: the strength of the links between its nodes. Two
questions turn on that structure; a third, the language-model failure signature, is taken
up in Section~\ref{sec:llm}.

\subsection{Adjudicating mechanism conflict}

When two stabilisers disagree, say perception reports one thing and well-attested
testimony another, a single-mechanism theory has no internal verdict to give. Reliabilism
can only ask which source is the more reliable here, which is the question at issue. The
network view adds a resource the component theories lack: the strength of the links
between the nodes, how tightly the mechanisms have co-varied, is itself evidence. Where
the stabilisers are densely connected and have tracked one another, a lone dissenter is
probably in error; where the links are weak, no single output earns much trust.
% TODO: this escapes relabeling ONLY if the connection strength (degree of co-variation),
% not the bare reliability of one mechanism, carries the epistemic weight. Build a worked
% example (perception vs instrument vs testimony) where the co-variation history flips the
% verdict a reliabilist would give. Cf. NOTES.md, "Claude's read", candidate 1.

\subsection{Domain-clustered co-variation}

The network view also predicts a cross-domain pattern no one-mechanism theory can
generate: because each component theory posits a single node, it has nothing to say about
how connectivity varies from one region to another. The stabilisers should be densely
linked where they've been jointly disciplined (everyday perception, basic coordination)
and sparsely linked where they haven't (theoretical physics, ethics, large-number
mathematics). Persistent disagreement should cluster in the sparsely connected regions.

The prediction is non-circular only if connectivity is characterised without appeal to
whether the mechanisms agree, and it can't be read off the graph in a purpose-free way:
\textcite{onishi_serpico_2022} show the abstraction level and boundaries stay
interest-relative. The reply is to fix the purpose rather than disown it. Holding the
representational aim of Section~\ref{sec:projectibility} fixed, connectivity is
characterised by shared selection history, institutional age, and measurement
infrastructure, each measurable without checking whether the mechanisms agree, and those
measures are what predict where disagreement clusters.
% TODO (the real work): operationalise the three connectivity proxies (selection history,
% institutional age, measurement infrastructure) and show on a worked case that they
% predict the disagreement pattern. Cf. NOTES.md, "Claude's read", candidate 2;
% Section~\ref{sec:cluster}.

\section{Objections}\label{sec:objections}
% TODO: deflationism; scientific revision; remote truth; mathematical/logical truth;
% moral/fictional/legal truth; perceptual illusion; the relabeling worry (answered in
% Section~\ref{sec:adjudication}); positioning vs alethic pluralism (Lynch/Wright/Pedersen).

\subsection{Interest-relativity}

One objection cuts at the kind-structure itself. \textcite{onishi_serpico_2022} argue that
Khalidi's causal-network account inherits the interest-relativity \textcite{craver_2009}
found in homeostatic-mechanism talk: which causal graph you draw, at what level of
abstraction, and where you set its boundaries all answer to explanatory interest, so the
network never fixes a unique, interest-free carving of kinds. If that's right,
truth-tracking-as-a-kind is only as objective as our purposes.

The reply concedes the premise and denies the damage. The account never needed an
interest-free carving. It makes two claims, both purpose-relative and both objective: a
system missing the perceptual and verificational stabilisers shows the failure signature
of Section~\ref{sec:llm}, and connectivity tracks selection history, institutional age, and
measurement infrastructure (Section~\ref{sec:adjudication}). Each holds once the
representational purpose is fixed, and neither requires that the purpose be read off nature.
\textcite{onishi_serpico_2022} reach the same general verdict, that interest-relativity is
ineliminable, and Khalidi grants domain pluralism; this account makes the operative
interest explicit instead of disowning it.

\section{The LLM payoff: participation, not a verdict}\label{sec:llm}
% Promise (1) from Section~\ref{sec:intro} is discharged HERE: hallucination is the
% predicted failure signature of a system with testimony + coherence but lacking
% perception + verification. WARNING (NOTES.md, "Claude's read", candidate 3): a bare
% "graded profile" is NOT HPC-specific -- the sibling (vector-grounding-problem_response)
% already owns that move. The HPC-specific payoff is the failure SIGNATURE, not mere
% unevenness. Do not re-run the sibling's argument.
% Dissolve the binary. LLMs participate in some truth-stabilising mechanisms and not
% others; grounding/truth-tracking becomes an empirical PROFILE. Mechanism-by-mechanism:
% perception (absent text-only / partial multimodal), testimony (strong but indirect),
% coherence (strong local / weak long-range), verification (mostly absent w/o tools),
% feedback/selection (RLHF tracks helpfulness/safety/preference, not truth alone),
% practical success (weak chat-only / stronger for agents). Connect to Mahowald et al.
% formal vs functional competence. This is the contribution boundary vs the two siblings.
TODO.

\section{Conclusion}\label{sec:conclusion}
TODO.

\section*{Acknowledgements}
% TODO: Replace with the actual models used on this project and their current versions.
% Check model names: frontier models change frequently.
This paper was drafted with the assistance of large language models. All content has been reviewed and revised by the author, who takes full responsibility for the final text.

\newpage
\printbibliography

\end{document}
